\documentclass[11pt]{article}

\usepackage[T1]{fontenc}
\usepackage[utf8]{inputenc}
\usepackage{lmodern}
\usepackage{microtype}
\usepackage[a4paper,margin=27mm,headheight=15pt]{geometry}
\usepackage{amsmath,amssymb,amsthm,mathtools}
\usepackage{booktabs,array,longtable,tabularx}
\usepackage{enumitem}
\usepackage{xcolor}
\usepackage{listings}
\usepackage{fancyhdr}
\usepackage{xurl}
\usepackage[hidelinks]{hyperref}
\usepackage{aliascnt}
\usepackage[nameinlink,noabbrev]{cleveref}

\allowdisplaybreaks
\setlength{\parindent}{0pt}
\setlength{\parskip}{0.55em}
\setcounter{tocdepth}{2}
\setcounter{secnumdepth}{3}

\definecolor{linkblue}{RGB}{16,73,132}
\hypersetup{
  colorlinks=true,
  linkcolor=linkblue,
  citecolor=linkblue,
  urlcolor=linkblue,
  pdftitle={Exact Geometric Tails and q-Binomial Extrapolation for Dyadic Values of Rvachev's Up-Function},
  pdfauthor={OpenAI}
}

\pagestyle{fancy}
\fancyhf{}
\fancyhead[L]{Dyadic Rvachev extrapolation}
\fancyhead[R]{\thepage}
\renewcommand{\headrulewidth}{0.4pt}

\newtheorem{theorem}{Theorem}[section]

\newaliascnt{proposition}{theorem}
\newtheorem{proposition}[proposition]{Proposition}
\aliascntresetthe{proposition}

\newaliascnt{lemma}{theorem}
\newtheorem{lemma}[lemma]{Lemma}
\aliascntresetthe{lemma}

\newaliascnt{corollary}{theorem}
\newtheorem{corollary}[corollary]{Corollary}
\aliascntresetthe{corollary}

\theoremstyle{definition}
\newaliascnt{definition}{theorem}
\newtheorem{definition}[definition]{Definition}
\aliascntresetthe{definition}

\newaliascnt{algorithm}{theorem}
\newtheorem{algorithm}[algorithm]{Algorithm}
\aliascntresetthe{algorithm}

\theoremstyle{remark}
\newaliascnt{remark}{theorem}
\newtheorem{remark}[remark]{Remark}
\aliascntresetthe{remark}

\newaliascnt{example}{theorem}
\newtheorem{example}[example]{Example}
\aliascntresetthe{example}

\crefname{proposition}{proposition}{propositions}
\Crefname{proposition}{Proposition}{Propositions}
\crefname{lemma}{lemma}{lemmas}
\Crefname{lemma}{Lemma}{Lemmas}
\crefname{corollary}{corollary}{corollaries}
\Crefname{corollary}{Corollary}{Corollaries}
\crefname{definition}{definition}{definitions}
\Crefname{definition}{Definition}{Definitions}
\crefname{algorithm}{algorithm}{algorithms}
\Crefname{algorithm}{Algorithm}{Algorithms}
\crefname{remark}{remark}{remarks}
\Crefname{remark}{Remark}{Remarks}
\crefname{example}{example}{examples}
\Crefname{example}{Example}{Examples}

\newcommand{\R}{\mathbb R}
\newcommand{\Q}{\mathbb Q}
\newcommand{\N}{\mathbb N}
\newcommand{\Z}{\mathbb Z}
\newcommand{\dd}{\,\mathrm d}
\newcommand{\sinc}{\operatorname{sinc}}
\newcommand{\rv}{\operatorname{up}}
\newcommand{\Fabius}{\mathcal F}
\newcommand{\Global}{\mathcal G}
\newcommand{\FT}{\mathscr F}
\newcommand{\D}{\mathrm D}
\newcommand{\coeff}{[z^r]}
\newcommand{\qbinom}[3]{\genfrac{[}{]}{0pt}{}{#1}{#2}_{#3}}
\newcommand{\floor}[1]{\left\lfloor #1\right\rfloor}
\newcommand{\abs}[1]{\left|#1\right|}
\newcommand{\norm}[1]{\left\lVert#1\right\rVert}
\newcommand{\e}{\mathrm e}
\newcommand{\ii}{\mathrm i}

\lstdefinestyle{code}{
  basicstyle=\ttfamily\small,
  columns=fullflexible,
  frame=single,
  breaklines=true,
  showstringspaces=false,
  backgroundcolor=\color{black!3},
  xleftmargin=0.4em,
  xrightmargin=0.4em
}

\title{\bfseries Exact Geometric Tails and \(q\)-Binomial Extrapolation\\[0.25em]
for Dyadic Values of Rvachev's Up-Function}
\author{A self-contained derivation based on the \texttt{ProveIt} development}
\date{2 September 2026}

\begin{document}
\maketitle

\begin{abstract}
Let \(\rv_n\) be the inverse Fourier transform of the finite product
\[
  \prod_{k=0}^{n}\sinc\!\left(\frac{\pi t}{2^k}\right),
  \qquad \sinc z=\frac{\sin z}{z},
\]
and fix a dyadic rational \(x\in[-1,1]\).  Write \(s=s(x)\) for the least
nonnegative integer for which \(2^s x\in\Z\), and put
\(R=\floor{s/2}\).  We prove the exact, eventually finite representation
\[
 \rv_n(x)=\rv(x)+\sum_{r=1}^{R}B_r(x)4^{-rn},\qquad n\ge s.
\]
Thus the error is not merely asymptotic: after at most \(s\) initial levels it
is a finite sum of signed geometric progressions, with ratios
\(4^{-1},4^{-2},\ldots,4^{-R}\).  The proof combines the exact dyadic spline
cells with an inverse convolution operator on a finite polynomial jet and the
fact that the derivative jet of \(\rv\) terminates at dyadic points.

The constant term is recovered from any \(R+1\) consecutive rational samples
by the single formula
\[
 \boxed{
 \rv(x)=\frac{1}{(Q;Q)_R}
 \sum_{j=0}^{R}(-1)^{R-j}Q^{\binom{R-j+1}{2}}
 \qbinom{R}{j}{Q}\,\rv_{n+j}(x),
 \qquad Q=\frac14,\quad n\ge s.}
\]
Here \((a;Q)_m\) is the \(q\)-Pochhammer symbol and
\(\qbinom{R}{j}{Q}\) is the Gaussian \(q\)-binomial coefficient.  We also
give an exact Vandermonde/Lagrange method for recovering every \(B_r(x)\), a
Romberg-style triangular implementation, stability estimates, examples, and
fully rational Python verification code.
\end{abstract}

\begin{remark}[Status and provenance]
The linked \texttt{ProveIt} corpus supplies the finite-spline formula, the
head--tail convolution identity, the Bernoulli--Bell reciprocal-tail
coefficients, the universal \(q\)-Lagrange weights, and the global derivative
formula used below \cite{proveit-primary,proveit-frontiers}.  The exact
\emph{termination at a fixed dyadic point} is derived here by joining those
ingredients through a local polynomial-jet inversion argument.  This article
is a mathematical proof and computational cross-check, not a claim that the
new combined statement has already been formalized in Lean.
\end{remark}

\tableofcontents
\newpage

\section{The result in one page}

Let
\begin{equation}
 \widehat f(t)=\int_{\R}f(x)\e^{-2\pi\ii tx}\dd x
 \label{eq:two-pi-fourier}
\end{equation}
be the Fourier convention, and normalize Rvachev's function by
\begin{equation}
 \widehat\rv(t)=\prod_{k=0}^{\infty}
 \sinc\!\left(\frac{\pi t}{2^k}\right).
 \label{eq:up-product}
\end{equation}
For \(n\ge0\), define the user's finite-product approximant by
\begin{equation}
 \widehat{\rv_n}(t)=\prod_{k=0}^{n}
 \sinc\!\left(\frac{\pi t}{2^k}\right).
 \label{eq:user-prefix}
\end{equation}

For a dyadic \(x\), its \emph{dyadic depth} is
\begin{equation}
 s(x)=\min\{s\ge0:2^s x\in\Z\}.
 \label{eq:depth}
\end{equation}
Thus \(s(0)=s(1)=s(-1)=0\), while a nonintegral dyadic has a unique reduced
form \(x=a/2^s\) with \(a\) odd and \(s=s(x)\).

The final answer is as follows.

\begin{theorem}[Eventual exact geometric representation]
\label{thm:main-preview}
Fix a dyadic \(x\in[-1,1]\), let \(s=s(x)\), and set
\(R=\floor{s/2}\).  There are rational numbers
\(B_1(x),\ldots,B_R(x)\) such that
\begin{equation}
 \boxed{
 \rv_n(x)=\rv(x)+\sum_{r=1}^{R}B_r(x)Q^{rn},
 \qquad Q=\frac14,\quad n\ge s.}
 \label{eq:main-preview}
\end{equation}
The coefficients have the analytic form
\begin{equation}
 B_r(x)=(-1)^r\alpha_r4^{-r}\rv^{(2r)}(x),
 \label{eq:B-derivative-preview}
\end{equation}
where the universal rational constants \(\alpha_r\) are determined by
\begin{equation}
 \sum_{r=0}^{\infty}\alpha_r z^r
 =\prod_{\ell=1}^{\infty}
 \frac{1}{\sinc(2^{-\ell}\sqrt z)}.
 \label{eq:alpha-product-preview}
\end{equation}
For every nonintegral dyadic of depth \(s\ge2\), the last coefficient
\(B_R(x)\) is nonzero, so the degree \(R\) cannot in general be reduced.
\end{theorem}

Once \cref{eq:main-preview} is known, extraction is finite interpolation at
geometric nodes.  Define
\begin{equation}
 (a;Q)_m=\prod_{k=0}^{m-1}(1-aQ^k),
 \qquad
 \qbinom{m}{j}{Q}=\frac{(Q;Q)_m}{(Q;Q)_j(Q;Q)_{m-j}}.
 \label{eq:q-notation-preview}
\end{equation}
Then, for any starting index \(n\ge s\),
\begin{equation}
 \boxed{
 \rv(x)=\frac{1}{(Q;Q)_R}
 \sum_{j=0}^{R}(-1)^{R-j}Q^{\binom{R-j+1}{2}}
 \qbinom{R}{j}{Q}\,\rv_{n+j}(x),
 \qquad Q=\frac14.}
 \label{eq:single-formula-preview}
\end{equation}
This is the requested single \(q\)-binomial/\(q\)-Pochhammer formula.  It
requires exactly \(R+1=\floor{s/2}+1\) rational spline values.

The rest of the article proves these statements, explains why the identity is
exact rather than merely asymptotic, and gives a reproducible implementation.

\section{Finite sinc products as exact dyadic splines}

\subsection{An indexing dictionary}

It is convenient to let \(p_N\) denote the inverse transform of a prefix with
exactly \(N\) factors:
\begin{equation}
 \widehat p_N(t)=\prod_{k=0}^{N-1}
 \sinc\!\left(\frac{\pi t}{2^k}\right),\qquad N\ge1.
 \label{eq:pN-prefix}
\end{equation}
The user's indexing and this factor-count indexing are related by
\begin{equation}
 \boxed{\rv_n=p_{n+1}.}
 \label{eq:index-dictionary}
\end{equation}
Keeping this one-step shift visible prevents the most common power-of-four
error in the final coefficients.

For \(a>0\), let
\begin{equation}
 b_a(x)=\frac{1}{2a}\mathbf 1_{[-a,a]}(x).
 \label{eq:box}
\end{equation}
Since \(\widehat b_a(t)=\sinc(2\pi at)\), Fourier inversion gives
\begin{equation}
 p_N=b_{2^{-1}}*b_{2^{-2}}*\cdots*b_{2^{-N}}.
 \label{eq:box-convolution}
\end{equation}
Hence \(p_N\) is a centered box spline, supported on
\begin{equation}
 [-A_N,A_N],\qquad A_N=\sum_{j=1}^{N}2^{-j}=1-2^{-N}.
 \label{eq:support-radius}
\end{equation}
It is piecewise polynomial of degree at most \(N-1\) and belongs to
\(C^{N-2}(\R)\) when \(N\ge2\).

\subsection{The Thue--Morse truncated-power formula}

Write
\begin{equation}
 \tau_k=(-1)^{s_2(k)},
 \label{eq:tau}
\end{equation}
where \(s_2(k)\) is the sum of the binary digits of \(k\).

\begin{proposition}[Exact finite spline]
\label{prop:truncated-power}
For \(N\ge1\), away from the harmless degree-zero knot convention,
\begin{equation}
 p_N(y)=\frac{2^{N(N-1)/2}}{(N-1)!}
 \sum_{k=0}^{2^N-1}\tau_k
 \left(y+A_N-\frac{k}{2^{N-1}}\right)_+^{N-1}.
 \label{eq:truncated-power}
\end{equation}
Its knots are
\begin{equation}
 t_{N,k}=-1+(2k+1)2^{-N},\qquad 0\le k<2^N.
 \label{eq:knots}
\end{equation}
In particular, \(p_N(y)\in\Q\) whenever \(y\in\Q\).
\end{proposition}

\begin{proof}
Differentiate every box in \cref{eq:box-convolution}.  Since
\(
 \D b_a=(2a)^{-1}(\delta_{-a}-\delta_a),
\)
the distribution \(\D^Np_N\) is a signed sum of \(2^N\) point masses.  Binary
endpoint choices give the signs \(\tau_k\), the locations in
\cref{eq:knots}, and the common multiplier
\(
 \prod_{j=1}^{N}2^{j-1}=2^{N(N-1)/2}.
\)
Integrating \(N\) times from \(-\infty\) yields
\cref{eq:truncated-power}.  Rationality at rational arguments is then
immediate.
\end{proof}

This formula is one of the central exact-arithmetic ingredients developed in
the repository \cite{proveit-frontiers}.

\subsection{The exact self-similar tail}

Let
\begin{equation}
 h_N=2^{-N},\qquad \rho_h(y)=h^{-1}\rv(y/h).
 \label{eq:scaled-tail}
\end{equation}
The omitted random/Fourier tail is a scaled copy of the full up-function.  More
precisely,
\begin{equation}
 \boxed{\rv=p_N*\rho_{h_N}.}
 \label{eq:head-tail}
\end{equation}
Indeed, the quotient between \cref{eq:up-product} and
\cref{eq:pN-prefix} is
\[
 \prod_{k=N}^{\infty}\sinc\!\left(\frac{\pi t}{2^k}\right)
 =\widehat\rv(2^{-N}t)=\widehat{\rho_{h_N}}(t).
\]
The identity is exact for every \(N\); no limiting argument is involved.

\section{Why dyadic points become cell centers}

The knot grid in \cref{eq:knots} has spacing \(2^{1-N}\).  The missing-tail
kernel in \cref{eq:head-tail} has radius \(h_N=2^{-N}\), exactly half that
spacing.  This numerical coincidence is the geometric mechanism behind the
finite answer.

\begin{lemma}[Dyadic midpoint lemma]
\label{lem:midpoint}
Let \(x\in[-1,1]\) have dyadic depth \(s\).  If \(N\ge s+1\), then the nearest
knots of \(p_N\) on either side of \(x\) are
\begin{equation}
 x-h_N\quad\text{and}\quad x+h_N,
 \label{eq:nearest-knots}
\end{equation}
with the natural exterior-cell interpretation at \(x=\pm1\).  Consequently,
there is a polynomial \(P_{N,x}\) of degree at most \(N-1\) such that
\begin{equation}
 p_N(x+t)=P_{N,x}(t),\qquad -h_N<t<h_N.
 \label{eq:local-polynomial}
\end{equation}
\end{lemma}

\begin{proof}
Since \(N-1\ge s\),
\[
 M=2^{N-1}(x+1)\in\Z.
\]
Therefore
\[
 x=-1+2M2^{-N},
 \qquad
 x\pm2^{-N}=-1+(2M\pm1)2^{-N}.
\]
The two latter numbers are consecutive members of the odd grid
\cref{eq:knots}.  Thus \(x\) is their midpoint and \(p_N\) is a single
polynomial throughout the open interval between them.
\end{proof}

In the user's indexing \(N=n+1\), the sufficient threshold is simply
\begin{equation}
 \boxed{n\ge s(x).}
 \label{eq:threshold}
\end{equation}
Earlier levels can occasionally satisfy the final formula by accident or by
symmetry, but \cref{eq:threshold} is the uniform, explicit guarantee.

\section{The reciprocal tail and its rational coefficients}

\subsection{Angular normalization}

The inverse differential operator is cleanest in the angular Fourier
variable.  Put
\begin{equation}
 \phi(\xi)=\widehat\rv\!\left(\frac{\xi}{2\pi}\right)
 =\prod_{\nu=1}^{\infty}\sinc\!\left(\frac{\xi}{2^\nu}\right).
 \label{eq:angular-product}
\end{equation}
Define the reciprocal-tail function
\begin{equation}
 \Phi(z)=\frac{1}{\phi(\sqrt z)}
 =\prod_{\ell=1}^{\infty}\frac{1}{\sinc(2^{-\ell}\sqrt z)}
 =\sum_{r=0}^{\infty}\alpha_r z^r,
 \qquad \alpha_0=1.
 \label{eq:Phi-alpha}
\end{equation}
The nearest singularity is at \(z=4\pi^2\), so this is an ordinary analytic
Taylor series near the origin.  For our polynomial-operator use, only its
formal coefficients matter.

\subsection{Bernoulli--Bell form}

The classical identity
\begin{equation}
 \log\sinc w=-\sum_{k=1}^{\infty}
 \frac{\zeta(2k)}{k\pi^{2k}}w^{2k}
 \label{eq:log-sinc}
\end{equation}
gives
\begin{equation}
 \Phi(z)=\exp\!\left(\sum_{k=1}^{\infty}c_kz^k\right),
 \label{eq:Phi-cumulants}
\end{equation}
where
\begin{equation}
 c_k=\frac{\zeta(2k)}{k\pi^{2k}(4^k-1)}
 =\frac{(-1)^{k+1}2^{2k-1}B_{2k}}
 {k(2k)!(4^k-1)}\in\Q.
 \label{eq:ck}
\end{equation}
It follows that every \(\alpha_r\) is rational.  In terms of complete
exponential Bell polynomials,
\begin{equation}
 \alpha_r=\frac{1}{r!}
 Y_r(1!c_1,2!c_2,\ldots,r!c_r),
 \label{eq:alpha-bell}
\end{equation}
and the computationally convenient recurrence is
\begin{equation}
 r\alpha_r=\sum_{k=1}^{r}k c_k\alpha_{r-k}.
 \label{eq:alpha-recurrence}
\end{equation}
The first values are
\begin{align}
 \alpha_0&=1,&
 \alpha_1&=\frac1{18},&
 \alpha_2&=\frac{31}{16200},\notag\\
 \alpha_3&=\frac{2347}{42865200},&
 \alpha_4&=\frac{1904369}{1311675120000},&
 \alpha_5&=\frac{3309760579}{88561680752160000}.
 \label{eq:first-alphas}
\end{align}
These are the same universal coefficients that occur in the repository's
all-order truncation expansion \cite{proveit-frontiers}.

\begin{remark}[Why an asymptotic expansion is not enough]
Fourier factorization gives, for each fixed truncation order \(M\),
\[
 p_N=\sum_{r=0}^{M}(-1)^r\alpha_r4^{-rN}\rv^{(2r)}
      +O(4^{-(M+1)N})
\]
in smooth norms.  At a dyadic point the displayed derivatives eventually
vanish, but this fact alone does \emph{not} force the remainder to be zero.
The exact result requires the local polynomial argument in the next section.
That is the step which upgrades ``all orders'' to a finite identity.
\end{remark}

\section{Exact inversion on a finite polynomial jet}

\subsection{Averaging a polynomial by the tail kernel}

Let \(Z\) have density \(\rv\), so \(Z\in[-1,1]\) and its law is even.  For a
polynomial \(P\) and \(h>0\), define
\begin{equation}
 (\mathcal C_hP)(t)=\int_{-1}^{1}\rv(z)P(t-hz)\dd z.
 \label{eq:Ch}
\end{equation}
This is again a polynomial of the same or lower degree.

\begin{lemma}[Finite polynomial inverse]
\label{lem:polynomial-inverse}
On the vector space of polynomials of degree at most \(d\),
\begin{equation}
 \boxed{
 P=\Phi(-h^2\D^2)\,\mathcal C_hP
 =\sum_{r=0}^{\floor{d/2}}(-1)^r\alpha_rh^{2r}
   \D^{2r}(\mathcal C_hP).}
 \label{eq:polynomial-inverse}
\end{equation}
The equality is exact and purely finite.
\end{lemma}

\begin{proof}
Let \(\mu_m=\int z^m\rv(z)\dd z\).  Odd moments vanish.  Taylor's formula for
a polynomial gives
\begin{equation}
 \mathcal C_h
 =\sum_{j\ge0}\frac{\mu_{2j}}{(2j)!}h^{2j}\D^{2j},
 \label{eq:Ch-moment}
\end{equation}
where the sum truncates on every finite-degree polynomial.  Its scalar symbol
is
\[
 M(w)=\sum_{j\ge0}\frac{\mu_{2j}}{(2j)!}w^j
     =\phi(\ii\sqrt w).
\]
By \cref{eq:Phi-alpha},
\[
 \Phi(-w)=\frac{1}{\phi(\sqrt{-w})}
         =\frac{1}{\phi(\ii\sqrt w)}
         =\frac1{M(w)}.
\]
Thus the formal power series of \(\Phi(-h^2\D^2)\) and \(\mathcal C_h\) are
reciprocals.  Because \(\D\) is nilpotent on polynomials of degree at most
\(d\), their composition is exactly the identity after finitely many terms;
there is no analytic convergence issue.
\end{proof}

\subsection{The spline jet equals the averaged local polynomial jet}

\begin{lemma}[Jet transfer at a cell center]
\label{lem:jet-transfer}
Let \(x\) have dyadic depth \(s\), let \(N\ge s+1\), set \(h=h_N\), and let
\(P_{N,x}\) be the local polynomial from \cref{lem:midpoint}.  Define
\begin{equation}
 J_{N,x}(t)=\mathcal C_hP_{N,x}(t).
 \label{eq:J-polynomial}
\end{equation}
Then, for every \(0\le k\le N-1\),
\begin{equation}
 J_{N,x}^{(k)}(0)=\rv^{(k)}(x).
 \label{eq:jet-match}
\end{equation}
Equivalently,
\begin{equation}
 J_{N,x}(t)=\sum_{k=0}^{N-1}\frac{\rv^{(k)}(x)}{k!}t^k.
 \label{eq:J-Taylor}
\end{equation}
\end{lemma}

\begin{proof}
From \cref{eq:head-tail},
\(
 \rv=p_N*\rho_h.
\)
For \(k\le N-1\), the distributional derivative \(\D^kp_N\) is locally
integrable; point masses first appear at order \(N\).  Therefore
\begin{align*}
 \rv^{(k)}(x)
 &=\int_{-h}^{h}(\D^kp_N)(x-y)\rho_h(y)\dd y\\
 &=\int_{-1}^{1}P_{N,x}^{(k)}(-hz)\rv(z)\dd z
  =J_{N,x}^{(k)}(0).
\end{align*}
The possible endpoint values at \(x\pm h\) affect a set of measure zero.  Both
sides of \cref{eq:J-Taylor} are polynomials of degree at most \(N-1\) with the
same derivatives at zero, so they are identical.
\end{proof}

Combining the preceding lemmas gives the exact local identity which drives the
whole result.

\begin{theorem}[Exact cell-center deconvolution]
\label{thm:cell-deconvolution}
If \(x\) has dyadic depth \(s\) and \(N\ge s+1\), then
\begin{equation}
 \boxed{
 p_N(x)=\sum_{r=0}^{\floor{(N-1)/2}}
 (-1)^r\alpha_r4^{-rN}\rv^{(2r)}(x).}
 \label{eq:cell-deconvolution}
\end{equation}
\end{theorem}

\begin{proof}
Apply \cref{lem:polynomial-inverse} to \(P_{N,x}\), with degree at most
\(N-1\), and evaluate at zero.  Since \(P_{N,x}(0)=p_N(x)\), while
\cref{lem:jet-transfer} gives
\(
 (\mathcal C_hP_{N,x})^{(2r)}(0)=\rv^{(2r)}(x),
\)
we obtain
\[
 p_N(x)=\sum_{r=0}^{\floor{(N-1)/2}}
 (-1)^r\alpha_r h^{2r}\rv^{(2r)}(x).
\]
Now \(h=2^{-N}\), so \(h^{2r}=4^{-rN}\).
\end{proof}

\section{The derivative jet terminates at a dyadic point}

\subsection{The signed global extension}

Let
\begin{equation}
 \Global(y)=\sum_{b\ge0}\tau_b\rv(y-2b-1).
 \label{eq:global-extension}
\end{equation}
This sum is locally finite.  On \([2b,2b+2]\), it is simply
\begin{equation}
 \Global(y)=\tau_b\rv(y-2b-1).
 \label{eq:global-block}
\end{equation}
In particular,
\begin{equation}
 \Global(2m)=0\qquad(m\in\Z_{\ge0}),
 \label{eq:G-even-zero}
\end{equation}
because even integers are support boundaries.  On \([-1,1]\),
\begin{equation}
 \rv(x)=\Global(x+1).
 \label{eq:up-G-translate}
\end{equation}
The global extension satisfies the dilation equation
\begin{equation}
 \Global'(y)=2\Global(2y).
 \label{eq:G-dilation}
\end{equation}
These identities are part of the arithmetic framework in
\cite{proveit-primary,arias2018}.

Iterating \cref{eq:G-dilation} gives
\begin{equation}
 \Global^{(m)}(y)=2^{\binom{m+1}{2}}\Global(2^m y),
 \label{eq:G-derivatives}
\end{equation}
and hence, from \cref{eq:up-G-translate},
\begin{equation}
 \boxed{
 \rv^{(m)}(x)=2^{\binom{m+1}{2}}
 \Global\!\left(2^m(x+1)\right),
 \qquad -1\le x\le1.}
 \label{eq:up-derivatives-G}
\end{equation}

\begin{lemma}[Finite dyadic jet]
\label{lem:dyadic-jet}
If \(x\) has dyadic depth \(s\), then
\begin{equation}
 \rv^{(m)}(x)=0\qquad\text{for every }m>s.
 \label{eq:dyadic-derivative-zero}
\end{equation}
\end{lemma}

\begin{proof}
For \(m>s\), the number \(2^m(x+1)\) is an even integer.  Indeed, if
\(s\ge1\), write \(x=a/2^s\) with \(a\) odd; then
\[
 2^m(x+1)=(a+2^s)2^{m-s},
\]
and \(m-s\ge1\).  The integral cases \(s=0\) are immediate as well.  Now use
\cref{eq:G-even-zero,eq:up-derivatives-G}.
\end{proof}

\begin{remark}[The last possible mode is genuinely present]
Suppose \(x=a/2^s\) is nonintegral and reduced.  If \(s=2R\), then
\(2^{2R}(x+1)=a+2^s\) is an odd integer, the center of a signed up-bump, so
\(\Global\) there is \(\pm1\).  If \(s=2R+1\), then
\(2^{2R}(x+1)=(a+2^s)/2\) is a half-integer strictly inside a signed bump, so
its \(\Global\)-value is nonzero.  Hence \(\rv^{(2R)}(x)\ne0\) whenever
\(s\ge2\).
\end{remark}

\section{Proof of the finite geometric formula}

We can now prove the main statement in its precise coefficient form.

\begin{theorem}[Exact finite geometric tail]
\label{thm:finite-geometric}
Let \(x\in[-1,1]\) be dyadic, let \(s=s(x)\), and put
\(R=\floor{s/2}\).  For every \(N\ge s+1\),
\begin{equation}
 \boxed{
 p_N(x)=\rv(x)+\sum_{r=1}^{R}C_r(x)4^{-rN},}
 \label{eq:pN-geometric}
\end{equation}
where
\begin{equation}
 C_r(x)=(-1)^r\alpha_r\rv^{(2r)}(x).
 \label{eq:C-r}
\end{equation}
Equivalently, in the user's indexing \(q_n=\rv_n(x)=p_{n+1}(x)\),
\begin{equation}
 \boxed{
 q_n=\rv(x)+\sum_{r=1}^{R}B_r(x)4^{-rn},
 \qquad n\ge s,}
 \label{eq:qn-geometric}
\end{equation}
with
\begin{align}
 B_r(x)
 &=4^{-r}C_r(x)
 =(-1)^r\alpha_r4^{-r}\rv^{(2r)}(x)
 \label{eq:B-r-derivative}\\
 &=(-1)^r\alpha_r2^{r(2r-1)}
 \Global\!\left(2^{2r}(x+1)\right).
 \label{eq:B-r-global}
\end{align}
All coefficients are rational.  If \(s\ge2\), then \(B_R(x)\ne0\).
\end{theorem}

\begin{proof}
The cell-center identity \cref{eq:cell-deconvolution} initially has terms up
to \(\floor{(N-1)/2}\).  By \cref{lem:dyadic-jet}, every derivative of order
larger than \(s\) is zero, so only \(2r\le s\), or \(r\le R\), remains.  This
proves \cref{eq:pN-geometric}.  Substitute \(N=n+1\) and absorb one factor
\(4^{-r}\) into the coefficient to obtain \cref{eq:qn-geometric}.
\Cref{eq:B-r-global} follows from
\cref{eq:up-derivatives-G}, because
\[
 4^{-r}2^{\binom{2r+1}{2}}
 =2^{-2r}2^{r(2r+1)}=2^{r(2r-1)}.
\]
The rationality of \(\alpha_r\) follows from \cref{eq:ck}.  Values of
\(\Global\) at dyadic arguments are rational by the standard arithmetic theory
of the Fabius/up function \cite{arias2018,proveit-primary}; alternatively,
rationality of every coefficient follows directly by solving the rational
Vandermonde system in \cref{sec:all-coefficients}.  The nonvanishing of the
last coefficient follows from the preceding remark and \(\alpha_R>0\).
\end{proof}

Thus the error sequence is a finite linear combination of signed geometric
progressions:
\begin{equation}
 q_n-\rv(x)=B_1(x)(1/4)^n+B_2(x)(1/16)^n+\cdots+B_R(x)(4^{-R})^n.
 \label{eq:signed-progressions}
\end{equation}
Each progression decreases in absolute value as \(n\) increases.  The number
of nonconstant modes is bounded by, and for \(s\ge2\) reaches, the explicit
order \(R=\floor{s/2}\).

\begin{corollary}[An exact recurrence]
\label{cor:recurrence}
Let \(E\) be the forward shift \((Eq)_n=q_{n+1}\).  On the tail \(n\ge s\),
\begin{equation}
 \prod_{r=1}^{R}(E-Q^r)\bigl(q_n-\rv(x)\bigr)=0,
 \qquad Q=\frac14,
 \label{eq:error-recurrence}
\end{equation}
and therefore
\begin{equation}
 (E-1)\prod_{r=1}^{R}(E-Q^r)q_n=0.
 \label{eq:q-recurrence}
\end{equation}
\end{corollary}

\begin{corollary}[Rational tail generating function]
For \(n\ge s\),
\begin{equation}
 \sum_{m=0}^{\infty}\bigl(q_{s+m}-\rv(x)\bigr)z^m
 =\sum_{r=1}^{R}\frac{B_r(x)Q^{rs}}{1-Q^rz}.
 \label{eq:tail-generating-function}
\end{equation}
\end{corollary}

\section{Exact extraction by \texorpdfstring{\(q\)}{q}-binomial interpolation}

\subsection{The interpolation problem}

Fix any starting index \(n\ge s\), set \(Q=1/4\), and rewrite
\cref{eq:qn-geometric} as
\begin{equation}
 q_{n+j}=L+\sum_{r=1}^{R}D_r(Q^j)^r,
 \qquad j=0,1,\ldots,R,
 \label{eq:scaled-interpolation-data}
\end{equation}
where
\begin{equation}
 L=\rv(x),\qquad D_r=B_r(x)Q^{rn}.
 \label{eq:D-r}
\end{equation}
Thus \(q_{n+j}=g_n(Q^j)\) for a polynomial
\begin{equation}
 g_n(z)=L+D_1z+\cdots+D_Rz^R.
 \label{eq:g-polynomial}
\end{equation}
The desired limit is simply \(L=g_n(0)\).

\subsection{Closed \texorpdfstring{\(q\)}{q}-Pochhammer weights}

For \(m\ge0\), recall
\begin{equation}
 (a;Q)_m=\prod_{k=0}^{m-1}(1-aQ^k),
 \qquad (Q;Q)_m=\prod_{k=1}^{m}(1-Q^k),
 \label{eq:q-Pochhammer}
\end{equation}
and
\begin{equation}
 \qbinom{m}{j}{Q}=\frac{(Q;Q)_m}
 {(Q;Q)_j(Q;Q)_{m-j}}.
 \label{eq:Gaussian-binomial}
\end{equation}

\begin{theorem}[Universal \(q\)-Lagrange weights]
\label{thm:q-weights}
For \(0<Q<1\), define
\begin{equation}
 w_j^{(R)}=
 \frac{(-1)^{R-j}Q^{\binom{R-j+1}{2}}}
 {(Q;Q)_j(Q;Q)_{R-j}}
 =\frac{(-1)^{R-j}Q^{\binom{R-j+1}{2}}}{(Q;Q)_R}
 \qbinom{R}{j}{Q}.
 \label{eq:q-weights}
\end{equation}
Then
\begin{equation}
 \sum_{j=0}^{R}w_j^{(R)}=1,
 \qquad
 \sum_{j=0}^{R}w_j^{(R)}Q^{jr}=0
 \quad(1\le r\le R).
 \label{eq:weight-moments}
\end{equation}
Consequently, for the dyadic spline sequence,
\begin{equation}
 \boxed{
 \rv(x)=\sum_{j=0}^{R}w_j^{(R)}\rv_{n+j}(x),
 \qquad n\ge s.}
 \label{eq:weighted-extraction}
\end{equation}
Equivalently,
\begin{equation}
 \boxed{
 \rv(x)=\frac{1}{(Q;Q)_R}
 \sum_{j=0}^{R}(-1)^{R-j}Q^{\binom{R-j+1}{2}}
 \qbinom{R}{j}{Q}\,\rv_{n+j}(x),
 \quad Q=\frac14,\ n\ge s.}
 \label{eq:main-extraction}
\end{equation}
\end{theorem}

\begin{proof}
The Lagrange coefficient for evaluating at zero from the nodes
\(1,Q,\ldots,Q^R\) is
\begin{equation}
 \ell_j(0)=\prod_{\substack{0\le k\le R\\k\ne j}}
 \frac{-Q^k}{Q^j-Q^k}.
 \label{eq:Lagrange-weight-product}
\end{equation}
Separate the factors with \(k<j\) and \(k>j\).  The signs, powers of \(Q\),
and products \(1-Q^m\) simplify to
\[
 \ell_j(0)=
 \frac{(-1)^{R-j}Q^{(R-j)(R-j+1)/2}}
 {(Q;Q)_j(Q;Q)_{R-j}},
\]
which is \cref{eq:q-weights}.  Lagrange interpolation is exact for every
polynomial of degree at most \(R\); applying it to \(1,z,\ldots,z^R\) yields
\cref{eq:weight-moments}.  Applying it to \(g_n\) in
\cref{eq:g-polynomial} yields \cref{eq:weighted-extraction}, and the Gaussian
form gives \cref{eq:main-extraction}.
\end{proof}

The first three filters, with \(Q=1/4\), are
\begin{align}
 R=1:\quad &\rv(x)=\frac{4q_{n+1}-q_n}{3},
 \label{eq:filter1}\\
 R=2:\quad &\rv(x)=\frac{q_n-20q_{n+1}+64q_{n+2}}{45},
 \label{eq:filter2}\\
 R=3:\quad &\rv(x)=
 \frac{-q_n+84q_{n+1}-1344q_{n+2}+4096q_{n+3}}{2835}.
 \label{eq:filter3}
\end{align}

\subsection{A triangular Richardson table}

The weighted sum can be evaluated without constructing all weights at once.

\begin{algorithm}[Exact \(q\)-Richardson table]
\label{alg:richardson}
Start with
\begin{equation}
 T_{0,n}=q_n.
 \label{eq:T0}
\end{equation}
For \(m=1,2,\ldots,R\), form
\begin{equation}
 \boxed{
 T_{m,n}=\frac{T_{m-1,n+1}-Q^mT_{m-1,n}}{1-Q^m}.}
 \label{eq:richardson-recursion}
\end{equation}
Then
\begin{equation}
 T_{R,n}=\rv(x)\qquad(n\ge s).
 \label{eq:richardson-exact}
\end{equation}
\end{algorithm}

At stage \(m\), the new row removes the mode \(Q^{mn}\).  This is the
Romberg-style realization of the polynomial
\begin{equation}
 W_R(E)=\prod_{m=1}^{R}\frac{E-Q^m}{1-Q^m}.
 \label{eq:filter-polynomial}
\end{equation}
The formulas and their moment identities are also developed in the repository's
\(q\)-Richardson layer \cite{proveit-frontiers}.

\section{Recovering every coefficient from the same samples}
\label{sec:all-coefficients}

The extraction formula is the first row of an ordinary rational Vandermonde
problem.  From \cref{eq:scaled-interpolation-data},
\begin{equation}
 \begin{pmatrix}
 1&1&1&\cdots&1\\
 1&Q&Q^2&\cdots&Q^R\\
 1&Q^2&Q^4&\cdots&Q^{2R}\\
 \vdots&\vdots&\vdots&&\vdots\\
 1&Q^R&Q^{2R}&\cdots&Q^{R^2}
 \end{pmatrix}
 \begin{pmatrix}L\\D_1\\D_2\\\vdots\\D_R\end{pmatrix}
 =
 \begin{pmatrix}q_n\\q_{n+1}\\q_{n+2}\\\vdots\\q_{n+R}\end{pmatrix}.
 \label{eq:Vandermonde-system}
\end{equation}
Its determinant is
\begin{equation}
 \prod_{0\le i<j\le R}(Q^j-Q^i)\ne0,
 \label{eq:Vandermonde-det}
\end{equation}
so the solution is unique.  Since the matrix and data are rational, every
\(L,D_r\), and hence every \(B_r=Q^{-rn}D_r\), is rational.

An explicit all-coefficient interpolation formula is
\begin{equation}
 g_n(z)=\sum_{j=0}^{R}q_{n+j}
 \prod_{\substack{0\le k\le R\\k\ne j}}
 \frac{z-Q^k}{Q^j-Q^k}.
 \label{eq:all-coefficient-Lagrange}
\end{equation}
Then
\begin{equation}
 \boxed{
 \rv(x)=g_n(0),\qquad
 B_r(x)=Q^{-rn}[z^r]g_n(z)\quad(1\le r\le R).}
 \label{eq:all-coefficients}
\end{equation}
Thus the same \(R+1\) values determine both the true rational limit and every
geometric correction coefficient.

\section{A practical exact-arithmetic method}

\begin{algorithm}[Limit extraction at a dyadic point]
\label{alg:practical}
Given a dyadic \(x\in[-1,1]\):
\begin{enumerate}[label=\arabic*.,leftmargin=2.2em]
 \item Reduce \(x\) and compute its dyadic depth
       \(s=\min\{j:2^jx\in\Z\}\).
 \item Set \(R=\floor{s/2}\), \(Q=1/4\), and choose any \(n\ge s\).  The
       minimal choice is \(n=s\).
 \item Evaluate the \(R+1\) rational spline values
       \(q_n,q_{n+1},\ldots,q_{n+R}\), for example by
       \cref{eq:truncated-power}.
 \item Apply the single formula \cref{eq:main-extraction}, or equivalently
       construct the triangular table \cref{eq:richardson-recursion}.  The
       result is exactly \(\rv(x)\).
 \item If the correction terms are wanted, solve
       \cref{eq:Vandermonde-system} or use
       \cref{eq:all-coefficients}.
 \item As an exact diagnostic, repeat the extraction from the next overlapping
       block.  Every block beginning at \(n\ge s\) must return the same
       rational number.
\end{enumerate}
\end{algorithm}

No estimate of a small remainder is needed.  With rational arithmetic, equality
is literal numerator/denominator equality.

\subsection{Stability for inexact inputs}

If the input samples are rounded rather than exact, the absolute condition
number of the weighted sum is
\begin{equation}
 \Lambda_R(Q)=\sum_{j=0}^{R}\abs{w_j^{(R)}}
 =\frac{(-Q;Q)_R}{(Q;Q)_R}.
 \label{eq:condition-number}
\end{equation}
For \(Q=1/4\), these values remain below the convergent limit
\begin{equation}
 \Lambda_\infty(1/4)=
 \frac{(-1/4;1/4)_\infty}{(1/4;1/4)_\infty}
 \approx1.9692603537.
 \label{eq:condition-limit}
\end{equation}
Thus even high-order dyadic filters have modest absolute amplification.  Exact
rational arithmetic, of course, eliminates roundoff entirely.

\subsection{Over-ordering}

If one uses any order \(M\ge R\), the same weight formula with \(M\) in place
of \(R\) also returns \(\rv(x)\) exactly, because it annihilates every actual
mode.  The minimal order \(R\) uses the fewest samples and normally gives the
smallest arithmetic workload.

\section{Worked exact examples}

\subsection{Depth two: one geometric mode}

For \(x=1/4\), \(s=2\) and \(R=1\).  Exact spline evaluation gives
\begin{equation}
 \rv_n\!\left(\frac14\right)
 =\frac{67}{72}+\frac19\,4^{-n},
 \qquad n\ge2.
 \label{eq:x-quarter-model}
\end{equation}
For example,
\[
 q_2=\frac{15}{16},\qquad q_3=\frac{179}{192},
\]
and \cref{eq:filter1} gives
\begin{equation}
 \rv\!\left(\frac14\right)
 =\frac{4(179/192)-15/16}{3}=\frac{67}{72}.
 \label{eq:x-quarter-extract}
\end{equation}
The fitted coefficient and \(\alpha_1=1/18\) imply
\(
 \rv''(1/4)=-8,
\)
consistent with \cref{eq:up-derivatives-G}.

Similarly,
\begin{equation}
 \rv_n\!\left(\frac34\right)
 =\frac{5}{72}-\frac19\,4^{-n},
 \qquad n\ge2,
 \label{eq:x-three-quarter-model}
\end{equation}
so the same filter returns \(\rv(3/4)=5/72\).

\subsection{Depth three: still one even-derivative mode}

For \(x=1/8\), \(s=3\) but \(R=\floor{3/2}=1\), because the correction
operator uses only even derivatives.  The exact identity is
\begin{equation}
 \rv_n\!\left(\frac18\right)
 =\frac{287}{288}+\frac1{18}\,4^{-n},
 \qquad n\ge3.
 \label{eq:x-eighth-model}
\end{equation}
Thus two consecutive values again suffice.

\subsection{Depth four: two geometric modes}

For \(x=1/16\), \(s=4\) and \(R=2\).  The three minimal samples are
\begin{equation}
 q_4=\frac{24575}{24576},\qquad
 q_5=\frac{1965959}{1966080},\qquad
 q_6=\frac{94365509}{94371840}.
 \label{eq:x-sixteenth-samples}
\end{equation}
The order-two filter gives
\begin{align}
 \rv\!\left(\frac1{16}\right)
 &=\frac{q_4-20q_5+64q_6}{45}
 =\frac{2073457}{2073600}.
 \label{eq:x-sixteenth-limit}
\end{align}
Solving the full system yields
\begin{equation}
 \boxed{
 \rv_n\!\left(\frac1{16}\right)
 =\frac{2073457}{2073600}
 +\frac{5}{648}4^{-n}
 -\frac{248}{2025}16^{-n},
 \qquad n\ge4.}
 \label{eq:x-sixteenth-model}
\end{equation}
The corresponding derivatives inferred from
\cref{eq:B-r-derivative} are
\begin{equation}
 \rv''\!\left(\frac1{16}\right)=-\frac59,
 \qquad
 \rv^{(4)}\!\left(\frac1{16}\right)=-1024.
 \label{eq:x-sixteenth-derivatives}
\end{equation}

\subsection{A table of exact models}

\begin{center}
\small
\setlength{\tabcolsep}{4pt}
\begin{tabular}{@{}c c c c c c@{}}
\toprule
\(x\) & \(s\) & \(R\) & \(\rv(x)\) & \(B_1\) & \(B_2\) \\
\midrule
\(0\) & 0 & 0 & \(1\) & -- & -- \\
\(1/2\) & 1 & 0 & \(1/2\) & -- & -- \\
\(1/4\) & 2 & 1 & \(67/72\) & \(1/9\) & -- \\
\(3/4\) & 2 & 1 & \(5/72\) & \(-1/9\) & -- \\
\(1/8\) & 3 & 1 & \(287/288\) & \(1/18\) & -- \\
\(1/16\) & 4 & 2 & \(2073457/2073600\) & \(5/648\) & \(-248/2025\) \\
\(3/16\) & 4 & 2 & \(2026943/2073600\) & \(67/648\) & \(248/2025\) \\
\(5/16\) & 4 & 2 & \(1767743/2073600\) & \(67/648\) & \(248/2025\) \\
\(1/32\) & 5 & 2 & \(33177581/33177600\) & \(1/2592\) & \(-124/2025\) \\
\bottomrule
\end{tabular}
\end{center}
In each row the model is
\(
 \rv_n(x)=\rv(x)+B_1 4^{-n}+B_2 16^{-n}
\)
with absent terms omitted and with validity from \(n\ge s\).

\section{Computational verification}

The archive accompanying this article contains
\begin{center}
 \texttt{exact\_extrapolation.py}
\end{center}
and its generated report
\begin{center}
 \texttt{verification\_output.txt}.
\end{center}
The program uses only the Python standard library.  Every quantity is a
\texttt{fractions.Fraction}; no floating-point value enters any identity check.
It implements:
\begin{itemize}[leftmargin=2em]
 \item the Thue--Morse truncated-power formula for \(p_N(x)\);
 \item \(q\)-Pochhammer symbols, Gaussian binomials, and the weights
       \cref{eq:q-weights};
 \item the triangular table \cref{eq:richardson-recursion};
 \item exact Gauss--Jordan solution of \cref{eq:Vandermonde-system};
 \item Bernoulli numbers, the cumulants \(c_k\), and the recurrence for
       \(\alpha_r\);
 \item overlapping-block and out-of-sample checks of every recovered model.
\end{itemize}

The included run verifies all signed dyadic points of depth at most six,
including endpoints: 129 points in total.  At each point it checks multiple
overlapping extraction blocks, the full fitted formula at later levels, and
the annihilating moment identities.  All checks pass exactly.  These
experiments are evidence against indexing or sign mistakes; the proof itself
is contained in \cref{lem:midpoint,lem:polynomial-inverse,lem:jet-transfer,lem:dyadic-jet,thm:finite-geometric,thm:q-weights}.

To reproduce the report from the archive, run
\begin{lstlisting}[style=code,language=bash]
python exact_extrapolation.py --max-depth 6 \
  --output verification_output.txt
\end{lstlisting}
Larger depths may be requested, at an exponentially growing cost for the
straightforward truncated-power evaluator.

\section{Further consequences and extensions}

\subsection{The universal part and the dyadic part}

Two logically distinct structures meet in the result:
\begin{enumerate}[label=(\roman*),leftmargin=2.2em]
 \item The reciprocal tail and the \(q\)-filter are universal for geometric
       sinc products.  For base \(b>1\), the natural ratio is \(Q=b^{-2}\),
       and the same weights \cref{eq:q-weights} annihilate the modes
       \(Q^n,\ldots,Q^{Rn}\).
 \item Exact finite termination here depends on special dyadic geometry:
       the point becomes the center of a uniform spline cell, and the global
       derivative orbit reaches an even integer after finitely many doublings.
\end{enumerate}
Thus \(q\)-Richardson acceleration exists much more generally, while exact
finite recovery is an arithmetic superconvergence phenomenon.

\subsection{Nonconsecutive levels}

Consecutive levels are not essential.  Given distinct levels
\(n_0,\ldots,n_R\ge s\), interpolate at the nodes
\(Q^{n_0},\ldots,Q^{n_R}\) and evaluate at zero:
\begin{equation}
 \rv(x)=\sum_{j=0}^{R}q_{n_j}
 \prod_{\substack{0\le k\le R\\k\ne j}}
 \frac{-Q^{n_k}}{Q^{n_j}-Q^{n_k}}.
 \label{eq:nonconsecutive}
\end{equation}
The consecutive choice is preferable because the common factor \(Q^n\)
cancels and produces the compact \(q\)-Pochhammer expression.

\subsection{Exact certification from data}

When \(x\) is known, the depth supplies a proof-level order bound, so there is
no model-selection problem.  Nevertheless, the following checks are useful in
software:
\begin{itemize}[leftmargin=2em]
 \item \(T_{R,n}=T_{R,n+1}\) for overlapping blocks;
 \item the recovered coefficients reproduce later \(q_k\) exactly;
 \item the moment sums in \cref{eq:weight-moments} vanish exactly;
 \item the highest fitted coefficient is nonzero for reduced depth \(s\ge2\).
\end{itemize}
Failure of any check signals a wrong index convention, an insufficient start
level, or an error in the spline evaluator.

\section{Conclusion}

At a fixed dyadic rational, the finite sinc-product splines exhibit a stronger
phenomenon than ordinary convergence.  From level \(n=s(x)\) onward, the point
is centered in one polynomial cell whose width matches the support of the
self-similar omitted tail.  The tail convolution can therefore be inverted
exactly on a finite polynomial jet.  The global doubling equation then kills
all derivatives above the dyadic depth.  What appears globally as an infinite
all-order expansion collapses pointwise to
\[
 \rv_n(x)=\rv(x)+\sum_{r=1}^{\floor{s(x)/2}}B_r(x)4^{-rn}.
\]

The constant term of this finite geometric polynomial is evaluation at zero
from the nodes \(1,Q,\ldots,Q^R\).  Lagrange interpolation on that geometric
grid is exactly the \(q\)-binomial filter
\[
 \rv(x)=\frac{1}{(Q;Q)_R}
 \sum_{j=0}^{R}(-1)^{R-j}Q^{\binom{R-j+1}{2}}
 \qbinom{R}{j}{Q}\,\rv_{n+j}(x),
 \qquad Q=\frac14.
\]
It uses finitely many rational inputs and returns the true rational value with
no limiting process and no remainder estimate.

\appendix

\section{Compact derivation of the weight formula}

For reference, start from
\[
 w_j=\prod_{k\ne j}\frac{-Q^k}{Q^j-Q^k}.
\]
For \(k<j\), write
\[
 Q^j-Q^k=-Q^k(1-Q^{j-k});
\]
for \(k>j\), write
\[
 Q^j-Q^k=Q^j(1-Q^{k-j}).
\]
There are \(j\) signs from the first group.  The denominator products are
\((Q;Q)_j\) and \((Q;Q)_{R-j}\).  The remaining power of \(Q\) is
\begin{align*}
 \sum_{k\ne j}k-
 \left(\sum_{k<j}k+j(R-j)\right)
 &=\frac{(R-j)(R-j+1)}2.
\end{align*}
Combining signs and powers gives \cref{eq:q-weights}.

\section{Normalization summary}

\begin{center}
\begin{tabularx}{\textwidth}{@{}>{\raggedright\arraybackslash}p{0.28\textwidth}X@{}}
\toprule
Object & Definition \\
\midrule
Full up-function &
\(\widehat\rv(t)=\prod_{k\ge0}\sinc(\pi t/2^k)\) under the two-pi convention.\\
User's approximant &
\(\widehat{\rv_n}(t)=\prod_{k=0}^{n}\sinc(\pi t/2^k)\).\\
Factor-count spline &
\(p_N=\rv_{N-1}\), containing exactly \(N\) factors.\\
Tail radius &
\(h_N=2^{-N}\), with \(\rv=p_N*\rho_{h_N}\).\\
Geometric ratio &
\(Q=1/4\).\\
Dyadic depth &
\(s(x)=\min\{s\ge0:2^sx\in\Z\}\).\\
Number of modes &
\(R=\floor{s(x)/2}\).\\
Guaranteed user level &
\(n\ge s(x)\).\\
\bottomrule
\end{tabularx}
\end{center}

\begingroup\sloppy\begin{thebibliography}{99}

\bibitem{proveit-primary}
V.~Reshetnikov,
\newblock \emph{The Fabius Function and Rvachev's Up-Function: Exact arithmetic,
probability, Fourier analysis, and corrected sharp asymptotics},
\newblock living TeX exposition in the \texttt{ProveIt} repository,
\newblock \url{https://github.com/VladimirReshetnikov/ProveIt/tree/main/Analysis/FabiusFunction/docs/Fabius_Function_and_Rvachev_Up}.

\bibitem{proveit-frontiers}
V.~Reshetnikov,
\newblock \emph{Non-formalized Research Frontiers: Exponents and q-Series},
\newblock living TeX development in the \texttt{ProveIt} repository,
\newblock \url{https://github.com/VladimirReshetnikov/ProveIt/tree/main/Analysis/FabiusFunction/docs/semi-formalized-research-frontiers}.

\bibitem{rvachev1990}
V.~A.~Rvachev,
\newblock Compactly supported solutions of functional-differential equations
and their applications,
\newblock \emph{Russian Mathematical Surveys} \textbf{45} (1990), no.~1,
87--120.

\bibitem{arias1982}
J.~Arias de Reyna,
\newblock Definici\'on y estudio de una funci\'on indefinidamente diferenciable
de soporte compacto,
\newblock \emph{Revista de la Real Academia de Ciencias Exactas, F\'isicas y
Naturales de Madrid} \textbf{76} (1982), no.~1, 21--38;
English translation, \emph{An infinitely differentiable function with compact
support: Definition and properties}, arXiv:1702.05442 (2017).

\bibitem{arias2018}
J.~Arias de Reyna,
\newblock Arithmetic of the Fabius function,
\newblock \emph{INTEGERS} \textbf{18} (2018), Article A51;
arXiv:1702.06487, version 3.

\bibitem{fabius1966}
J.~Fabius,
\newblock A probabilistic example of a nowhere analytic \(C^\infty\)-function,
\newblock \emph{Zeitschrift f\"ur Wahrscheinlichkeitstheorie und Verwandte
Gebiete} \textbf{5} (1966), 173--174.

\bibitem{deboor}
C.~de Boor,
\newblock \emph{A Practical Guide to Splines}, revised edition,
\newblock Springer, New York, 2001.

\bibitem{gasper-rahman}
G.~Gasper and M.~Rahman,
\newblock \emph{Basic Hypergeometric Series}, second edition,
\newblock Cambridge University Press, 2004.

\end{thebibliography}\endgroup

\end{document}
